% Example cover page for a two-page IEEE conference document
%\onecloumn
% Cover page

\begin{center}
  \textbf{recovery testbed}
\end{center}

\begin{itemize}
  \item DEADLINE: Abstract April 11, Paper April 18, 2025
  \item Publication Venue: QEST+FORMATS 2025
  \item Url for Publication Venue: \url{https://www.qest-formats.org/}
  \item Submission Category or Track:
    \begin{itemize}
        \item Models and metrics for the correctness, performance, reliability, safety, and security of systems (stochastic, probabilistic, and non-deterministic models including Markov chains, automata, Petri nets, process algebra, max-plus algebra, and their variations);
    \end{itemize}
  \item Page limit:  16 pages; currently at 16
  \item Research Contribution: 
   The proposed model enables comparisons of alternative
repair schemes by considering the availability of repair crews,
predicting failure propagation, and estimating the level of
service capacity over time. The model can also guide staffing
decisions by identifying the minimum number of repair crews
required to maintain a particular service level.
\end{itemize}

\vspace{2ex}

%SSS- This helps us keep track of the number of drafts
\textbf{Version 1.1}
Last Updated: 4/1/2025 \\
\textbf{Major Changes Since Previous Draft}
\begin{itemize}
    \item Case study
    \item Conclusion
\end{itemize}

\vspace{2ex}

\textbf{Outline}
\begin{enumerate}
%SSS - Include the number of pages in your current draft in front of each section (not subsection) title. Also include the status (complete, pending, not written).
  %SSS - The following sections should be in every paper. 
  \item  Abstract
  \item Introduction 
  \begin{enumerate}
          \item Blackout in Kenya
          \item Motivation for this paper
          \item Intro to this model
        \end{enumerate}
  \item Related work
  \begin{enumerate}
          \item Modeling of cascading failure
          \item Modeling of the Recovery process
        \end{enumerate}
  \item Proposed recovery model
  \begin{enumerate}
          \item Modeling of Cascading failure
          \begin{enumerate}
              \item States
              \item Actions
              \item Rewards
              \item Transition probabilities
          \end{enumerate}
          \item Repair schemes
          \begin{enumerate}
              \item Randomized priorities
              \item pre-made repair priorities
              \item Look-ahead dynamic priorities
              \item cascading effect dynamic priorities
          \end{enumerate}
        \end{enumerate}
  \item case studies
  \begin{enumerate}
          \item IEEE-14
          \item Simulation Environment
          \item IEEE-14 Failure Cases
          \begin{enumerate}
            \item experiment 1
            \item experiment 2
            \item Experiment 3
            \end{enumerate}
        \item Discussion of results
        \end{enumerate}
  \item Conclusion and Future Work 
    \begin{enumerate}
          \item The contribution to the body of knowledge 
          \item social impact 
          \item future directions, enable advances in the repair scheme area
        \end{enumerate}
  \item Acknowledgements 
  %SSS - Include this text in your acknowledgments: The authors gratefully acknowledge funding from NSF DUE-1742523, DEd GAANN awards P200A180095 and P200A210121, and the Missouri S\&T Intelligent Systems Center.
  \item References
\end{enumerate}

\vspace{2ex}

\textbf{TO DO (Advisors):}
\begin{itemize}
    \item review the paper
\end{itemize}

\textbf{TO DO (Sanfan):}
\begin{itemize}
    \item review the paper entirely
    \item submit for publication 
\end{itemize}

\textbf{Author guidelines (submission requirements)}: \\ 
\url{https://www.qest-formats.org/cfp.html}\\
Regular papers must not exceed 16 pages, excluding references.
We are considering additional presentation-only submission types; details will be provided in the later CfPs.

All papers must be submitted in Springer’s LNCS format and will undergo a rigorous single-blind review process. All submitted papers must be unpublished and not be submitted for publication elsewhere. Papers should be submitted electronically using EasyChair at this link.

Springer encourages authors to include their ORCIDs in their papers. Authors should consult Springer’s authors’ guidelines and use Springer’s LaTeX templates for the preparation of their papers. Submitted papers not complying with the above guidelines may be rejected without undergoing review.
We are considering awards for best papers and best artifacts.

Publications\\
All accepted papers need to be presented and discussed at the conference by one of the authors. The QEST+FORMATS 2025 proceedings will be published in the Springer LNCS series indexed by ISI Web of Science, Scopus, ACM Digital Library, dblp, Google Scholar. All submitted papers will be evaluated by at least three reviewers on the basis of their originality, technical quality, scientific or practical contribution to the state of the art, methodology, clarity, and adequacy of references.

Special Issue\\
We are working on a special issue in a Q1 journal. A selection of the best papers will be invited to submit an extended version of their work to special issues in internationally recognized journals.

Artifact Evaluation\\
Reproducibility of experimental results is crucial to foster an atmosphere of trustworthy, open, and reusable research. To improve and reward reproducibility, QEST+FORMATS 2025 will include a dedicated Artifact Evaluation (AE). Submission of an artifact is mandatory for tool papers (both regular and short), and optional but encouraged for research and case study papers where it can support the results presented in the paper. More details can be found in the corresponding page
%End Cover Sheet

